\section{System Requirements (Requisiti di Sistema)}

I requisiti di sistema formalizzano in dettaglio le capacità funzionali, i vincoli qualitativi e regole derivanti dal dominio, che la piattaforma software \textbf{MyAma} deve soddisfare.

% =========================================================================
% 4.1 REQUISITI FUNZIONALI
% =========================================================================
\subsection{Requisiti Funzionali}
Questa sezione descrive le funzioni che il sistema deve fornire, strutturate in base agli attori che interagiscono con la piattaforma \textbf{MyAma}, derivate dall'analisi dei casi d'uso.

\begin{center}
\renewcommand{\arraystretch}{1.3}
\begin{longtable}{|p{3.0cm}|p{1.2cm}|p{4.0cm}|p{5.8cm}|}
\caption{Requisiti Funzionali}\label{tab:rf} \\ \hline
\rowcolor{tableheader} \textbf{Attore} & \textbf{ID} & \textbf{Requisito Funzionale} & \textbf{Descrizione / Criterio di Verifica} \\ \hline
\endfirsthead

\hline
\rowcolor{tableheader} \textbf{Attore} & \textbf{ID} & \textbf{Requisito Funzionale} & \textbf{Descrizione / Criterio di Verifica} \\ \hline
\endhead

Utente non registrato & \textbf{RF-01} & Registrazione cittadino & Il sistema deve consentire a un utente di registrarsi come cittadino fornendo i propri dati e accettando l'informativa privacy. \\ \hline

Utente non registrato & \textbf{RF-02} & Registrazione tramite invito & Il sistema deve permettere la registrazione del personale (Autista, Operatore o Amministratore) tramite l'inserimento di un codice di invito valido. \\ \hline

Utente di sistema & \textbf{RF-03} & Autenticazione & Il sistema deve permettere agli utenti registrati di effettuare l'accesso inserendo credenziali valide, reindirizzandoli alle funzionalità del proprio ruolo. \\ \hline

Cittadino & \textbf{RF-04} & Richiesta ritiro a domicilio & Il sistema deve permettere al cittadino di prenotare un ritiro a domicilio fornendo indirizzo, CAP e i dettagli del rifiuto ingombrante. \\ \hline

Cittadino & \textbf{RF-05} & Fornire sedi coperte & Il sistema deve fornire all'utente le sedi AMA coperte dal servizio richiesto in base al CAP o la zona specificata dal Cittadino nella configurazione del suo profilo. \\ \hline

Cittadino & \textbf{RF-06} & Prenotazione conferimento & Il sistema deve permettere al cittadino di prenotare il conferimento del rifiuto presso una sede AMA compatibile con la propria zona. \\ \hline

Cittadino & \textbf{RF-07} & Annullamento prenotazione & Il sistema deve consentire al cittadino di annullare una propria prenotazione precedentemente creata. \\ \hline

Autista AMA & \textbf{RF-08} & Visualizzazione ritiri assegnati & Il sistema deve fornire all'autista l'elenco dei ritiri a domicilio assegnati per il proprio turno, con i dettagli dell'indirizzo e del rifiuto. \\ \hline

Autista AMA & \textbf{RF-09} & Registrazione esito ritiro & Il sistema deve permettere all'autista di registrare l'esito (es. positivo o negativo) del ritiro a domicilio effettuato. \\ \hline

Operatore di Sede & \textbf{RF-10} & Verifica e registrazione conferimento & Il sistema deve permettere all'operatore di verificare una prenotazione in sede e registrarne l'esito. \\ \hline

Amministratore Sede & \textbf{RF-11} & Gestione risorse e disponibilità & Il sistema deve permettere all'amministratore di configurare le fasce orarie e la disponibilità di veicoli e lavoratori. \\ \hline

Amministratore Sede & \textbf{RF-12} & Generazione codici invito & Il sistema deve consentire all'amministratore di generare codici di invito per registrare e abilitare nuovo personale. \\ \hline

Amministratore Sede & \textbf{RF-13} & Rimozione personale dalla sede & Il sistema deve permettere all'amministratore di sede di rimuovere e disassociare membri del personale operativo dalla propria struttura. \\ \hline

Amministratore Generale & \textbf{RF-14} & Generazione codici invito amministratori & Il sistema deve consentire all'amministratore generale di generare codici di invito abilitanti al ruolo di amministratore di sede. \\ \hline

Amministratore Generale & \textbf{RF-15} & Revoca e rimozione amministratori di sede & Il sistema deve consentire all'amministratore generale di revocare e rimuovere l'abilitazione degli amministratori di sede AMA. \\ \hline
\end{longtable}
\end{center}

% =========================================================================
% 4.2 REQUISITI NON FUNZIONALI
% =========================================================================
\subsection{Requisiti Non Funzionali}
Questa sezione descrive i vincoli di qualità e le caratteristiche misurabili che il sistema deve rispettare.

\begin{center}
\renewcommand{\arraystretch}{1.3}
\begin{longtable}{|p{1.8cm}|p{3.2cm}|p{9.0cm}|}
\caption{Requisiti Non Funzionali}\label{tab:rnf} \\ \hline
\rowcolor{tableheader} \textbf{ID} & \textbf{Requisito} & \textbf{Descrizione e Metriche (Criterio di Verifica)} \\ \hline
\endfirsthead

\hline
\rowcolor{tableheader} \textbf{ID} & \textbf{Requisito} & \textbf{Descrizione e Metriche (Criterio di Verifica)} \\ \hline
\endhead

\textbf{RNF-01} & Performance & Le operazioni di registrazione e rimozione degli account, annullamento delle prenotazioni e registrazione dell'esito di ritiri e conferimenti devono completarsi entro \textbf{2 secondi} in condizioni operative normali. Le operazioni di prenotazione di un servizio e di generazione dei codici di invito devono completarsi entro \textbf{3 secondi}. \\ \hline

\textbf{RNF-02} & Scalabilità & Il sistema deve poter gestire l'aumento del numero di cittadini, personale AMA, sedi e prenotazioni senza richiedere modifiche sostanziali all'architettura e senza compromettere i requisiti di performance definiti. L'incremento del carico deve poter essere gestito mediante l'aggiunta di risorse elaborative. \\ \hline

\textbf{RNF-03} & Portabilità & Il sistema deve essere accessibile tramite i principali browser moderni --- Chrome, Firefox, Safari ed Edge --- e adattare l'interfaccia alle diverse dimensioni dello schermo mediante responsive design. L'accesso deve essere possibile sia da dispositivi desktop sia mobili. \\ \hline

\textbf{RNF-04} & Affidabilità & Il sistema deve preservare l'integrità e la consistenza dei dati anche in caso di anomalie o interruzioni improvvise. Operazioni critiche quali creazione, annullamento e aggiornamento dello stato delle prenotazioni non devono produrre stati parziali o incoerenti. In caso di errore l'operazione deve essere completata correttamente oppure annullata senza alterare i dati preesistenti. \\ \hline

\textbf{RNF-05} & Manutenibilità & Il sistema deve essere organizzato secondo una struttura modulare e disaccoppiata, accompagnata da documentazione tecnica sufficiente a permettere la modifica o sostituzione dei singoli componenti limitando l'impatto sulle altre parti del sistema. \\ \hline

\textbf{RNF-06} & Disponibilità & Il sistema deve garantire una disponibilità del servizio pari almeno al \textbf{99,9\%}, escluse le finestre di manutenzione programmata. Le funzionalità rivolte ai cittadini devono essere accessibili \textbf{24 ore su 24 e 7 giorni su 7}, mentre quelle destinate al personale AMA devono essere disponibili durante gli orari operativi previsti. Gli interventi di manutenzione programmata devono essere effettuati preferibilmente nelle fasce di minor utilizzo. In caso di guasto critico il servizio deve poter essere ripristinato entro \textbf{3 ore}. \\ \hline

\textbf{RNF-07} & Usabilità & Il sistema deve offrire un'interfaccia semplice e intuitiva. Un cittadino deve poter completare una prenotazione in un massimo di \textbf{5 passaggi guidati}, con controlli che prevengano o segnalino gli errori di inserimento. Le funzionalità destinate al personale AMA devono essere organizzate in modo da ridurre al minimo la formazione necessaria. \\ \hline

\textbf{RNF-08} & Sicurezza & Il sistema deve garantire riservatezza e integrità dei dati personali e aziendali. Le comunicazioni di rete devono utilizzare protocolli cifrati e le password devono essere memorizzate mediante algoritmi di hashing sicuri. L'accesso alle funzionalità e ai dati deve essere regolato mediante controllo degli accessi basato sui ruoli (\textbf{RBAC}): ciascun utente può accedere esclusivamente alle operazioni previste dal proprio ruolo e, per il personale AMA, ai dati necessari allo svolgimento delle attività della propria sede o competenza.\newline\newline L'autenticazione dell'utente dovrebbe avvenire mediante l'uso di token di autenticazione per validare la sessione, che il server può rilasciare all'utente al momento del login. Successivamente l'utente specifica al server il token ad ogni richiesta. \\ \hline
\end{longtable}
\end{center}

% =========================================================================
% 4.3 REQUISITI DI DOMINIO
% =========================================================================
\subsection{Requisiti di Dominio}
Questa sezione raccoglie le regole aziendali o logiche specifiche del contesto ``AMA'' (smaltimento rifiuti) che vincolano il funzionamento del software:

\begin{center}
\renewcommand{\arraystretch}{1.3}
\begin{longtable}{|p{1.8cm}|p{3.2cm}|p{9.0cm}|}
\caption{Requisiti di Dominio}\label{tab:rd} \\ \hline
\rowcolor{tableheader} \textbf{ID} & \textbf{Regola di Dominio} & \textbf{Descrizione} \\ \hline
\endfirsthead

\hline
\rowcolor{tableheader} \textbf{ID} & \textbf{Regola di Dominio} & \textbf{Descrizione} \\ \hline
\endhead

\textbf{RD-01} & Competenza Territoriale & Una richiesta di ritiro a domicilio o la scelta di una sede per il conferimento sono vincolate dalla competenza territoriale e dalle fasce orarie operative disponibili delle sedi. Il sistema può accettare richieste solo se il CAP o la zona indicata nella richiesta utente sono coperti e assegnati a una Sede AMA attiva. \\ \hline

\textbf{RD-02} & Vincolo di Capacità dei Veicoli & Durante la pianificazione dei ritiri a domicilio, il sistema non può assegnare a un singolo veicolo un numero di prenotazioni il cui peso o volume stimato superi il carico massimo consentito (capacità) del veicolo stesso per quel turno. Inoltre il rifiuto è ritirabile solo se conforme a quanto dichiarato, e se la sede o il veicolo assegnato sono abilitati a gestirlo. \\ \hline

\textbf{RD-03} & Coerenza del Ciclo di Vita delle Prenotazioni & Il sistema deve garantire la coerenza del ciclo di vita della prenotazione attraverso gli stati ammessi (\textit{In attesa} $\rightarrow$ \textit{Confermata} $\rightarrow$ \textit{In corso} $\rightarrow$ \textit{Completata}, oppure \textit{Annullata} / \textit{Non eseguita}), impedendo transizioni di stato illecite quali l'annullamento o la modifica di prenotazioni già in lavorazione o concluse. \\ \hline

\textbf{RD-04} & Validità del Codice di Invito & Un codice di invito generato da un amministratore di sede o generale deve essere univoco, associato a un ruolo specifico e deve essere monouso con scadenza temporale, per evitare la creazione fraudolenta di profili con privilegi elevati. \\ \hline

\textbf{RD-05} & Controllo degli Accessi basato sul Ruolo (RBAC) & Il sistema deve garantire che ogni utente possa accedere esclusivamente alle informazioni e alle funzionalità previste dal proprio ruolo. I dati personali del Cittadino, come informazioni anagrafiche, indirizzo e dati relativi alle prenotazioni, devono essere accessibili solo al personale AMA autorizzato e quando necessario allo svolgimento del servizio. Analogamente, gli Amministratori di sede devono poter operare solamente sui dati e sulle risorse delle sedi di propria competenza, mentre l'Amministratore Generale può accedere alle informazioni necessarie alla gestione complessiva del sistema. \\ \hline
\end{longtable}
\end{center}
